% RESEARCH_DRAFT_SHA256: e0c5f893dbc947409533058019a2d9ea9d2ce5004e2e3f36d62bef29818199f3
% Layout and source-reference wrappers expanded for the web renderer.
\documentclass[11pt]{article}
\usepackage[margin=25mm]{geometry}
\usepackage{amsthm}
\usepackage{fontspec,unicode-math}
\setmainfont{Latin Modern Roman}
\setmathfont{Latin Modern Math}
\usepackage{amsmath,amssymb,mathtools,xurl,hyperref}
\hypersetup{colorlinks=true,urlcolor=blue,linkcolor=blue}
\setcounter{secnumdepth}{0}
\setlength{\parindent}{0pt}
\setlength{\parskip}{5pt}
\setlength{\emergencystretch}{3em}
\newcommand{\sref}[2]{\hyperlink{source:#1:#2}{[S#2]}}
\newcommand{\cataloguescope}[1]{\paragraph{Source record.}#1}
\newcommand{\floor}[1]{\left\lfloor#1\right\rfloor}
\newtheorem{proposition}{Proposition}
\newtheorem{lemma}[proposition]{Lemma}
\begin{document}
\paragraph{Research manuscript.} OpenAI \texttt{gpt-6-astra}, reasoning \texttt{ultra}; 11 October 2026.
\paragraph{Preserved source body.} The complete catalogue mathematical text and sources follow with layout and source links expanded, including source qualifications and its historical preparation notes. New research begins at the section headed Research attempt.

% UnsolvedMath ID 30003867; source catalogue code OWR-16170-001
\section{Statistical Recovery of High-Dimensional Network Structure}\label{30003867}
{\small UnsolvedMath ID 30003867 · Computer Science}
\subsection{Definitions and mathematical statement}

\paragraph{Shared notation.}Standard number systems are $\mathbb N,\mathbb Z,\mathbb Q,\mathbb R,\mathbb C$. Specialized notation and index ranges follow the cited source and the definitions below.


A network model specifies a probability distribution for its edges or observations. Parameter estimation fits that distribution to data; recovering a structure also requires a criterion distinguishing true from incorrect network descriptions.
\paragraph{Record scope and qualifications.}
The catalogue title is retained as metadata. Named certificates, contracts, bounds or frameworks in that title are not asserted as results of the original source extract.
The network distribution, sample access, and estimation target are unspecified.

Source-record limitation: The supplied assessment flags an underspecified network-estimation prompt with no model or guarantee.
\paragraph{Statement recorded in the supplied source.}
Given observations from a high-dimensional network model, how can the underlying network structure, including its edge-generation rule, be statistically estimated? \href{https://doi.org/10.4171/owr/2018/29}{[S2]}
\paragraph{Source status.}
The supplied catalogue labels this source record open; that metadata was checked against the supplied snapshot on 2026-10-11. The label is preserved for provenance and is not an independent verification of present open status or mathematical correctness. The catalogue literature assessment is dated 2026-08-21.
\subsection{Short English statement}
Estimate high-dimensional network structure and the rule generating its edges.
\subsection{Sources}
\begin{itemize}
\item[S1] Ulam AI and the UnsolvedMath Contributors, UnsolvedMath: A Curated Collection of Open Mathematics Problems, record 30003867 (OWR-16170-001). Catalogue curation and metadata: CC BY 4.0. The source extract is retained; definitions, English summaries and qualifications are editorial additions.
\url{https://huggingface.co/datasets/ulamai/UnsolvedMath}.
\item[S2] Matrix Estimation Meets Statistical Network Analysis: Extracting low-dimensional structures in high dimension (2019), Oberwolfach Reports, DOI 10.4171/owr/2018/29, pp. 1745--1783
\url{https://doi.org/10.4171/owr/2018/29}.
\item[S3] Oberwolfach Report 29/2018
\url{https://ems.press/content/serial-article-files/46753}.
\item Source-specific Oberwolfach Reports licences, including CC BY-SA 4.0 where applicable, are retained; see THIRD\_PARTY\_NOTICES.md.
\end{itemize}
\paragraph{Preparation and review.}
The original source and UnsolvedMath contributors retain attribution for the source extract. Codex (GPT-6; exact serving revision not exposed) assisted with definitions, English summaries, scope notes and formatting. No new human mathematical authorship or independently verified solution is asserted. Permanent source ID 30003867 and code OWR-16170-001 checked against the complete supplied registry and current repository. Distinct upstream records, including repeated questions and broad or damaged extracts, are retained for complete coverage of the source CS category. A repeated record is not a new independent conjecture. Complete mathematical-equivalence and current-literature review has not been performed. Preparation and final review follow contribution prompts C and D, with source gaps disclosed. Structural and rendering checks do not establish a new theorem or open status.
\label{end:30003867}

\section{Research attempt}
\subsection{Literature, source context and scope}
The full original catalogue body and source qualifications are preserved above. The relevant primary workshop passages were read in the cached publisher report. Historical claims and subsequent attributed results are distinguished from the mathematical work below. A restricted example, standard lemma or known theorem is not counted as new progress on the entire original target. Missing assumptions and repeated or rhetorical extracts are not silently turned into different solved problems.

\subsection{A graph-generating rule is identifiable only with a convention}
The source's introductory network-estimation discussion explicitly warns that the word network is ambiguous. We test recovery of a graph-generating kernel in a specified latent-position model. Draw independent $U_i\sim\operatorname{Unif}[0,1]$ and, conditional on them, include edges independently with probabilities $W(U_i,U_j)$. Compare $W(x,y)=xy$ with $\widetilde W(x,y)=(1-x)(1-y)$.
\begin{proposition}
These two kernels generate exactly the same distribution of every finite observed graph, although $\|W-\widetilde W\|_{L^2([0,1]^2)}^2=1/6$. No graph-only estimator can, for both kernels, have probability greater than $1/2$ of an $L^2$ error smaller than $1/(2\sqrt6)$ relative to their displayed coordinates.
\end{proposition}
\begin{proof}
In the second model couple its latent positions as $\widetilde U_i=1-U_i$. They are still independent uniform variables and satisfy $\widetilde W(\widetilde U_i,\widetilde U_j)=U_iU_j$. Using the same independent edge uniforms makes the observed graphs identical. Further,
\[
W(x,y)-\widetilde W(x,y)=x+y-1,
\]
whose squared integral is the variance of the sum of two independent uniform variables, namely $1/12+1/12=1/6$. Equivalently this follows by integrating the expanded quadratic. The two open $L^2$ balls of radius $1/(2\sqrt6)$ are disjoint by the triangle inequality. Coupling the estimator's random seed gives identical estimated kernels in the two worlds, so the success indicators sum to at most one. Their probabilities cannot both exceed $1/2$.
\end{proof}
\subsection{Verification and remaining gaps}
The quadratic integral was evaluated exactly by monomial integration, and the coupled edge probabilities were checked on a rational grid. The failed route is estimating a particular coordinate representation before accounting for latent relabeling. The two kernels are equivalent under a measure-preserving coordinate change, so this does not obstruct estimation in the natural quotient by that symmetry. No rates, estimator for that quotient, or conclusion about every high-dimensional network model are established. The source's broad model-selection agenda is not replaced by this auxiliary graphon problem.

\section{References and source attribution}
Primary workshop report: \url{https://doi.org/10.4171/owr/2018/29}. Relevant passages read: 2018/29 p. 1746, network-estimation framing and ambiguity in edge-generation rules.

All retrievals were checked on 11 October 2026. Original question attribution remains in the preserved Sources section. External theorems are attributed and not claimed as new discoveries or fully verified merely by citation.


\paragraph{Integration review.} Codex (GPT-6; exact serving revision not exposed) reviewed the written arguments and expanded layout and source-reference wrappers. Repository metadata and the production web renderer were checked using Python 3.13 and JavaScript with MathJax 3.2.2. This model-assisted review is not independent human proof verification; no PDF compilation or formal proof checking is asserted.

\section{Completion Estimate}
\noindent\textbf{COMPLETION: 0\%}
No substantiated new progress toward the entire original target is claimed. The worked results and failed route are delimited above; a known result, illustrative model, or auxiliary obstruction is not counted as original-target completion.

\section{Final status and attribution}
\textbf{Attempt status: unresolved.} This describes this research attempt, not a new ruling on the source catalogue's status or on external literature claims.
The mathematics and this self-audit were generated by OpenAI \texttt{gpt-6-astra}, reasoning \texttt{ultra}, on 11 October 2026. Codex (GPT-6; exact serving revision not exposed) orchestrated the batch, reviewed the written mathematical arguments, and contributed formatting/repository review. This additional model review is not independent human verification. No human mathematical authorship, independent human proof checking, or formal verification is claimed. The historical preparation paragraphs in the reproduced catalogue body concern that body only.

Contribution prompt A was applied to the complete definition and supplied source record, followed by the adversarial checks of prompt B. Local reproduction files are \texttt{batch022/checks.py}, \texttt{computations.json}, and the per-ID audit JSON. They were not public links at preparation time. Numerical checks are finite and do not prove unrestricted results. No LaTeX engine was found locally; compilation was not performed by this research agent.

\end{document}
